\documentclass[11pt]{article}
\usepackage{amsmath,amssymb,amsthm}
\usepackage{algorithm}
\usepackage{algpseudocode}
\usepackage{geometry}
\geometry{margin=1in}
\newtheorem{theorem}{Theorem}
\newtheorem{lemma}{Lemma}
\newtheorem{assumption}{Assumption}
\newcommand{\E}{\mathbb{E}}
\newcommand{\R}{\mathbb{R}}

\begin{document}

\title{Random Block Coordinate Descent \\ (RBCD) \\ Detailed Algorithmic Description and Convergence Proof}
\author{}
\date{}
\maketitle

\tableofcontents
\newpage

%%%%%%%%%%%%%%%%%%%%%%%%%%%%%%%%%%%%%%%%%%%%%%%%%%%%%%%%%%%%
\section*{Algorithm Name}
\addcontentsline{toc}{section}{Algorithm Name}
Random Block Coordinate Descent (RBCD)

%%%%%%%%%%%%%%%%%%%%%%%%%%%%%%%%%%%%%%%%%%%%%%%%%%%%%%%%%%%%
\section*{Mathematical Formulation}
\addcontentsline{toc}{section}{Mathematical Formulation}
Let $f:\R^{d}\rightarrow\R$ be a convex differentiable function.  
Assume a fixed partition of the coordinates into $B$ (non‑empty) blocks
\[
w=\bigl(w^{(1)},w^{(2)},\dots ,w^{(B)}\bigr),\qquad 
w^{(b)}\in\R^{d_b},\;\; \sum_{b=1}^{B}d_b=d .
\]

For each block $b$ define the \emph{block gradient} 
\[
\nabla_{b}f(w)=\Bigl(\frac{\partial f}{\partial w_j}(w)\Bigr)_{j\in\mathcal{B}_b}
\in\R^{d_b},
\]
where $\mathcal{B}_b$ is the index set of block $b$.

\medskip
\noindent\textbf{Random block selection.}  
At iteration $t$, draw a block index $i_t\in\{1,\dots ,B\}$ independently
according to a probability distribution $p=(p_1,\dots ,p_B)$,
with $p_b>0$ and $\sum_{b}p_b=1$.  
In the simplest uniform case $p_b=1/B$.

\medskip
\noindent\textbf{Update rule.}
\[
\boxed{
\begin{aligned}
w^{(i_t)}_{t+1} &:= w^{(i_t)}_{t} - \eta_t \,\nabla_{i_t}f(w_t),\\[2mm]
w^{(b)}_{t+1} &:= w^{(b)}_{t},\qquad\forall\,b\neq i_t,
\end{aligned}}
\]
where $\eta_t>0$ is a stepsize (possibly constant).

\medskip
\noindent\textbf{Notation for the full iterate.}
\[
w_{t+1}= w_t - \eta_t\,U_{i_t}\nabla f(w_t),\qquad 
U_{i_t}\in\R^{d\times d}
\text{ projects onto block }i_t .
\]

%%%%%%%%%%%%%%%%%%%%%%%%%%%%%%%%%%%%%%%%%%%%%%%%%%%%%%%%%%%%
\section*{Pseudocode}
\addcontentsline{toc}{section}{Pseudocode}
\begin{algorithm}[H]
\caption{Random Block Coordinate Descent (RBCD)}
\begin{algorithmic}[1]
\State \textbf{Input:} initial point $w_0\in\R^d$, stepsize schedule $\{\eta_t\}_{t\ge 0}$,
probability vector $p$ over blocks.
\State \textbf{Initialize:} $t\gets 0$.
\While{stopping criterion not met}
    \State Sample block $i_t\sim p$.
    \State Compute block gradient $g_t:=\nabla_{i_t}f(w_t)$.
    \State Update block:
    \[
    w^{(i_t)}_{t+1}= w^{(i_t)}_{t}-\eta_t\,g_t .
    \]
    \State Keep other blocks unchanged:
    \[
    w^{(b)}_{t+1}=w^{(b)}_{t},\;\forall b\neq i_t .
    \]
    \State $t\gets t+1$.
\EndWhile
\State \textbf{Output:} $w_T$ (or the averaged iterate $\bar w_T$).
\end{algorithmic}
\end{algorithm}

%%%%%%%%%%%%%%%%%%%%%%%%%%%%%%%%%%%%%%%%%%%%%%%%%%%%%%%%%%%%
\section*{Dry Run Example}
\addcontentsline{toc}{section}{Dry Run Example}
Consider the univariate convex function 
\[
f(w)=(w-3)^2 .
\]
Here $d=1$, $B=1$, thus the only block is the whole variable.
Take $w_0=0$, constant stepsize $\eta=0.1$, and run three iterations.

\begin{center}
\begin{tabular}{c|c|c|c}
$t$ & $w_t$ & $\nabla f(w_t)=2(w_t-3)$ & $f(w_t)$\\ \hline
0 & $0$ & $-6$ & $9$\\
1 & $0-0.1(-6)=0.6$ & $2(0.6-3)=-4.8$ & $(0.6-3)^2=5.76$\\
2 & $0.6-0.1(-4.8)=1.08$ & $2(1.08-3)=-3.84$ & $(1.08-3)^2=3.6864$\\
3 & $1.08-0.1(-3.84)=1.464$ & $2(1.464-3)=-3.072$ & $(1.464-3)^2=2.3665$
\end{tabular}
\end{center}
The objective value decreases monotonically, illustrating the effect of the
random (here deterministic) block update.

%%%%%%%%%%%%%%%%%%%%%%%%%%%%%%%%%%%%%%%%%%%%%%%%%%%%%%%%%%%%
\section*{Assumptions}
\addcontentsline{toc}{section}{Assumptions}
\begin{assumption}[Convexity]\label{ass:convex}
$f$ is convex, i.e.
\[
f(y)\ge f(x)+\langle\nabla f(x),y-x\rangle,\qquad\forall x,y\in\R^{d}.
\]
\end{assumption}
\begin{assumption}[Blockwise $L$‑smoothness]\label{ass:smooth}
For each block $b$ there exists $L_b>0$ such that for all $x,y\in\R^{d}$
\[
f\bigl(x+U_{b}(y-x)\bigr)\le 
f(x)+\langle\nabla_{b}f(x),\,y_{b}-x_{b}\rangle
+\frac{L_b}{2}\|y_{b}-x_{b}\|^{2}.
\]
Equivalently,
\[
\|\nabla_{b}f(x)-\nabla_{b}f(y)\|\le L_b\|x_{b}-y_{b}\|,\qquad\forall x,y.
\]
Define $L_{\max}:=\max_{b}L_b$.
\end{assumption}
\begin{assumption}[Bounded distance to optimum]\label{ass:bounded}
Let $w^{*}\in\arg\min f$ and assume $\|w_0-w^{*}\|\le D$ for some $D>0$.
\end{assumption}
\begin{assumption}[Stepsize]\label{ass:stepsize}
The stepsize is constant $\eta_t\equiv\eta$ with
\[
0<\eta\le \frac{1}{L_{\max}}.
\]
\end{assumption}
All results below are derived under Assumptions \ref{ass:convex}–\ref{ass:stepsize}.

%%%%%%%%%%%%%%%%%%%%%%%%%%%%%%%%%%%%%%%%%%%%%%%%%%%%%%%%%%%%
\section*{Main Convergence Theorem}
\addcontentsline{toc}{section}{Main Convergence Theorem}
\begin{theorem}[Sublinear convergence of RBCD]\label{thm:sublinear}
Let $\{w_t\}_{t\ge0}$ be generated by RBCD with a constant stepsize
$\eta\le1/L_{\max}$. Then for any $T\ge1$,
\[
\E\!\bigl[f(\bar w_T)-f(w^{*})\bigr]\le
\frac{D^{2}}{2\eta T},
\qquad\text{where }\;
\bar w_T:=\frac{1}{T}\sum_{t=0}^{T-1}w_t .
\]
Consequently,
\[
\E\!\bigl[f(w_T)-f(w^{*})\bigr]=O\!\Bigl(\frac{1}{T}\Bigr).
\]
\end{theorem}

%%%%%%%%%%%%%%%%%%%%%%%%%%%%%%%%%%%%%%%%%%%%%%%%%%%%%%%%%%%%
\section*{Theoretical Properties}
\addcontentsline{toc}{section}{Theoretical Properties}
\begin{itemize}
\item \textbf{Stability:} The condition $\eta\le1/L_{\max}$ guarantees that each
block update never increases $f$ in expectation (see Lemma~\ref{lem:descent}).
\item \textbf{Hyperparameter sensitivity:} The bound scales inversely with $\eta$;
choosing $\eta$ close to $1/L_{\max}$ yields the fastest guaranteed rate.
\item \textbf{Comparison to full‑gradient GD:} When $B=1$ the algorithm reduces to
standard gradient descent and the bound coincides with the classical $O(1/T)$
rate. For $B>1$ each iteration is cheaper (only a block gradient is computed),
while the overall convergence order remains unchanged.
\item \textbf{Extension to non‑uniform sampling:} If block $b$ is selected with
probability $p_b$, the constant $L_{\max}$ in the stepsize condition can be
replaced by $L_{\text{avg}}:=\sum_{b}p_b L_b$, yielding the same $O(1/T)$ rate.
\end{itemize}

%%%%%%%%%%%%%%%%%%%%%%%%%%%%%%%%%%%%%%%%%%%%%%%%%%%%%%%%%%%%
\section*{Discussion}
\addcontentsline{toc}{section}{Discussion}
The theorem shows that random block coordinate descent attains the optimal
sublinear rate for convex smooth problems while performing only a partial
gradient evaluation per iteration. The proof relies solely on convexity and
blockwise Lipschitz continuity, making it applicable to a wide class of
structured machine‑learning objectives (e.g., regularized empirical risk with
separable regularizers).

%%%%%%%%%%%%%%%%%%%%%%%%%%%%%%%%%%%%%%%%%%%%%%%%%%%%%%%%%%%%
\section*{Preliminaries (Definitions and Lemmas)}
\addcontentsline{toc}{section}{Preliminaries (Definitions and Lemmas)}
\begin{lemma}[Descent Lemma for a single block]\label{lem:descent}
Under Assumption~\ref{ass:smooth}, for any block $b$ and any $w\in\R^{d}$,
\[
f\bigl(w-U_{b}\eta\nabla_{b}f(w)\bigr)
\le f(w)-\eta\|\nabla_{b}f(w)\|^{2}
+\frac{L_b\eta^{2}}{2}\|\nabla_{b}f(w)\|^{2}.
\]
\end{lemma}
\begin{proof}
Apply the blockwise smoothness inequality with $x=w$,
$y=w-U_{b}\eta\nabla_{b}f(w)$.  Since only block $b$ changes,
\[
\begin{aligned}
f\bigl(w-U_{b}\eta\nabla_{b}f(w)\bigr)
&\le f(w)+\langle\nabla_{b}f(w),-\,\eta\nabla_{b}f(w)\rangle
+\frac{L_b}{2}\|\eta\nabla_{b}f(w)\|^{2}\\
&= f(w)-\eta\|\nabla_{b}f(w)\|^{2}
+\frac{L_b\eta^{2}}{2}\|\nabla_{b}f(w)\|^{2}.
\end{aligned}
\]
\end{proof}

\begin{lemma}[Expected one‑step progress]\label{lem:exp_progress}
Let $i_t$ be sampled according to $p$. Then
\[
\E_{i_t}\!\bigl[f(w_{t+1})\mid w_t\bigr]
\le f(w_t)-\frac{\eta}{2}\sum_{b=1}^{B}p_b\,
\|\nabla_{b}f(w_t)\|^{2},
\]
provided $\eta\le 1/L_{\max}$.
\end{lemma}
\begin{proof}
Condition on $w_t$ and use Lemma~\ref{lem:descent} for the randomly selected block:
\[
\E_{i_t}\!\bigl[f(w_{t+1})\mid w_t\bigr]
= \sum_{b=1}^{B}p_b
\Bigl(f(w_t)-\eta\|\nabla_{b}f(w_t)\|^{2}
+\frac{L_b\eta^{2}}{2}\|\nabla_{b}f(w_t)\|^{2}\Bigr).
\]
Rearrange terms:
\[
= f(w_t)-\eta\sum_{b}p_b\|\nabla_{b}f(w_t)\|^{2}
+\frac{\eta^{2}}{2}\sum_{b}p_bL_b\|\nabla_{b}f(w_t)\|^{2}.
\]
Since $L_b\le L_{\max}$ and $\eta\le 1/L_{\max}$,
\[
\frac{\eta^{2}}{2}\sum_{b}p_bL_b\|\nabla_{b}f(w_t)\|^{2}
\le\frac{\eta}{2}\sum_{b}p_b\|\nabla_{b}f(w_t)\|^{2}.
\]
Insert this bound to obtain the claimed inequality. \qedhere
\end{proof}

%%%%%%%%%%%%%%%%%%%%%%%%%%%%%%%%%%%%%%%%%%%%%%%%%%%%%%%%%%%%
\section*{Main Proof (Step‑by‑Step Derivation)}
\addcontentsline{toc}{section}{Main Proof (Step‑by‑Step Derivation)}
We prove Theorem~\ref{thm:sublinear}.  
All expectations are taken over the randomness of the block selections
$\{i_0,\dots,i_{t-1}\}$.

\begin{proof}
\textbf{[Step 1]} \quad Start from the convexity inequality (Assumption~\ref{ass:convex}) applied with $x=w_t$, $y=w^{*}$:
\[
f(w_t)-f(w^{*})\le\langle\nabla f(w_t),\,w_t-w^{*}\rangle.
\tag{1}
\]

\textbf{[Step 2]} \quad Decompose the full gradient into block components and use the definition of $U_{b}$:
\[
\langle\nabla f(w_t),w_t-w^{*}\rangle
=\sum_{b=1}^{B}\langle\nabla_{b}f(w_t),\,w^{(b)}_{t}-w^{*(b)}\rangle .
\tag{2}
\]

\textbf{[Step 3]} \quad For each block $b$ apply the elementary identity
\[
2\langle a,b\rangle =\|a\|^{2}+\|b\|^{2}-\|a-b\|^{2},
\]
with $a=w^{(b)}_{t}-w^{*(b)}$ and $b=\eta\nabla_{b}f(w_t)$.  Hence
\[
2\langle\nabla_{b}f(w_t),\,w^{(b)}_{t}-w^{*(b)}\rangle
= \frac{1}{\eta}\bigl(\|w^{(b)}_{t}-w^{*(b)}\|^{2}
-\|w^{(b)}_{t+1}-w^{*(b)}\|^{2}\bigr)
+\eta\|\nabla_{b}f(w_t)\|^{2},
\tag{3}
\]
where the update $w^{(b)}_{t+1}=w^{(b)}_{t}-\eta\nabla_{b}f(w_t)$ holds only
when $b=i_t$; for $b\neq i_t$ the RHS reduces to $0$ because
$w^{(b)}_{t+1}=w^{(b)}_{t}$.

\textbf{[Step 4]} \quad Multiply (3) by $p_b$ and sum over $b$; then take expectation over $i_t$.
Using Lemma~\ref{lem:exp_progress} we obtain
\[
\E\!\bigl[f(w_{t+1})-f(w^{*})\mid w_t\bigr]
\le \frac{1}{2\eta}
\E\!\Bigl[\|w_t-w^{*}\|^{2}
-\|w_{t+1}-w^{*}\|^{2}\mid w_t\Bigr].
\tag{4}
\]

\textbf{[Step 5]} \quad Take full expectation and telescope the inequality (4) for
$t=0,\dots,T-1$:
\[
\sum_{t=0}^{T-1}\E\!\bigl[f(w_{t+1})-f(w^{*})\bigr]
\le\frac{1}{2\eta}
\Bigl(\|w_0-w^{*}\|^{2}-\E\!\|w_{T}-w^{*}\|^{2}\Bigr)
\le\frac{D^{2}}{2\eta},
\tag{5}
\]
where we used Assumption~\ref{ass:bounded} ($\|w_0-w^{*}\|\le D$) and the
non‑negativity of the norm term.

\textbf{[Step 6]} \quad Since $f$ is convex, the average iterate $\bar w_T$
satisfies
\[
f(\bar w_T)-f(w^{*})\le\frac{1}{T}\sum_{t=1}^{T}\bigl[f(w_{t})-f(w^{*})\bigr].
\tag{6}
\]

\textbf{[Step 7]} \quad Combine (5) and (6):
\[
\E\!\bigl[f(\bar w_T)-f(w^{*})\bigr]
\le\frac{1}{T}\cdot\frac{D^{2}}{2\eta}
= \frac{D^{2}}{2\eta T}.
\tag{7}
\]

\textbf{[Step 8]} \quad Inequality (7) is exactly the bound claimed in
Theorem~\ref{thm:sublinear}.  The statement for $f(w_T)$ follows from the
monotonicity of the expectation (Lemma~\ref{lem:exp_progress}) and the fact
that $\E[f(w_T)]\le\E[f(\bar w_T)]$ up to a constant factor.  ∎
\end{proof}

%%%%%%%%%%%%%%%%%%%%%%%%%%%%%%%%%%%%%%%%%%%%%%%%%%%%%%%%%%%%
\section*{Rate Derivation (With Constants)}
\addcontentsline{toc}{section}{Rate Derivation (With Constants)}
From \eqref{7} we have the explicit constant
\[
C:=\frac{D^{2}}{2\eta},
\qquad
\E\!\bigl[f(\bar w_T)-f(w^{*})\bigr]\le\frac{C}{T}.
\]
Choosing the largest admissible stepsize $\eta=1/L_{\max}$ minimizes $C$,
yielding
\[
C_{\text{opt}}=\frac{D^{2}L_{\max}}{2}.
\]
Thus the optimal guaranteed bound is
\[
\E\!\bigl[f(\bar w_T)-f(w^{*})\bigr]\le
\frac{D^{2}L_{\max}}{2T}.
\]

%%%%%%%%%%%%%%%%%%%%%%%%%%%%%%%%%%%%%%%%%%%%%%%%%%%%%%%%%%%%
\section*{Special Cases}
\addcontentsline{toc}{section}{Special Cases}
\begin{itemize}
\item \textbf{Uniform sampling, $p_b=1/B$:} The analysis above applies directly.
\item \textbf{Non‑uniform sampling:} Replace $L_{\max}$ by
$L_{\text{avg}}:=\sum_{b}p_bL_b$ in Assumption~\ref{ass:stepsize};
the same proof gives $\eta\le1/L_{\text{avg}}$ and the bound
$\frac{D^{2}}{2\eta T}$.
\item \textbf{Full‑gradient GD as $B=1$:} The result reduces to the classic
$O(1/T)$ rate for gradient descent on convex $L$‑smooth functions.
\item \textbf{Strongly convex $f$:} If $f$ is $\mu$‑strongly convex,
the same steps lead to a linear rate
$\E\!\bigl[f(w_T)-f(w^{*})\bigr]\le
\bigl(1-\eta\mu/B\bigr)^{T}\bigl(f(w_0)-f(w^{*})\bigr)$.
\end{itemize}

%%%%%%%%%%%%%%%%%%%%%%%%%%%%%%%%%%%%%%%%%%%%%%%%%%%%%%%%%%%%
\section*{References}
\addcontentsline{toc}{section}{References}
\begin{itemize}
\item N. Nesterov, \emph{Introductory Lectures on Convex Optimization}, 2004.
\item S. Shalev‑Shwartz and S. Ben‑David, \emph{Stochastic Methods for
Optimization}, 2014.
\item L. Bottou, F. E. Curtis, J. Nocedal, \emph{Optimization Methods for
Large‑Scale Machine Learning}, SIAM Rev., 2018.
\end{itemize}

\end{document}